\MWChapter{Conclusion and next decisions}{结论与后续决策}
\label{ch:conclusion}

Cantilune proposes a single formal object for visible, evolvable orchestration: typed composition
in \(\mathcal C\), identified reconfiguration in \(\mathcal R\), and certified readings into
dependency, resource, communication, and compositional views. The current repository evidence
binds a generic Core Theory to a verified source and build record; that is a meaningful formal
milestone, but it is not a runtime or product launch \citep{cantiluneReadme,cantiluneClosure}.

\begin{MWDecision}
\textbf{Current decision: retain pre-alpha research status.} Publish the formal boundary and
evidence chain as \emph{proved / review-pending}; do not describe Cantilune as deployed,
independently reviewed, FCP-approved, or product-conformant until those distinct records exist.
\end{MWDecision}

\section{Prioritized next work / 优先后续工作}

\begin{enumerate}
  \item Obtain independent category/DPO/Petri, \(\pi\)/domain, and Lean/provenance reviews with
    explicit conflict-of-interest records.
  \item Complete the applicable FCP and ADR decision path without widening the proof claims beyond
    their declared D1-A and projection boundaries.
  \item Define one concrete product package and provide its complete per-rule conformance
    certificates, including admission, resource, authorization, fairness, and replay evidence.
  \item Only after a runtime exists, prepare a separate implementation report for interfaces,
    security, usability, performance, reliability, and operational safety.
\end{enumerate}

\section{Reader checklist / 读者检查表}

When reusing or extending this report, preserve the distinction between specification,
mechanization, evidence binding, independent review, and product conformance. Update the source
branch and commit identifiers before copying a theorem status, and remove any empirical or
deployment language that is not supported by the target project's own artifacts.
